\section{Primer eje central: entrenar el modelo para que sea un Agent nativo}

\subsection{La conciencia del problema del agentic training}

Varios invitados situaron ``native agentic'' como el problema de investigación central de 2025. La formulación de Pan Jiayi (潘嘉怡) es particularmente precisa: este es el año en que la comunidad empezó a pensar de manera sistemática en ``cómo hacer que un modelo se vuelva native agentic''. La palabra clave de este problema no es solo ``saber usar herramientas'', sino si el modelo, después del entrenamiento, puede tener capacidades más estables de coding, computer use, multi-agent collaboration y reasoning.

\textit{``Este es el primer año en que la gente empieza en serio a pensar de manera sistemática en cómo hacer que un modelo se vuelva native agentic.''}

\begin{importantbox}{El objetivo del agentic training no es agregar unas cuantas etiquetas de habilidad más}
Lo que en realidad se busca resolver es la estructura de capacidades del modelo:
\begin{itemize}
  \item si, ante una tarea abierta, puede descomponerla y explorarla por iniciativa propia;
  \item si, ante la retroalimentación del entorno, puede corregir su acción;
  \item si, ante tareas de cadena larga, puede mantener el estado intermedio;
  \item si, ante llamadas a herramientas de múltiples rondas, puede formar una estrategia estable.
\end{itemize}
\end{importantbox}

\subsection{El escenario de data science ilustra el valor del entrenamiento}

El ``deep analysis'' que hizo el equipo de Zhang Shaolei (张少磊) ofrece una explicación de aplicación muy contundente. La ciencia de datos no es una tarea de conversación pura: exige que el modelo alterne entre bases de datos, scripts, métricas de evaluación y semántica de negocio. Con solo un workflow de capa externa se puede armar el sistema en un principio, pero si el modelo mismo no tiene buenas capacidades de análisis, codificación e interacción con el entorno, el sistema pierde estabilidad muy pronto en tareas de cadena larga. Por eso inyectaron aún más, mediante aprendizaje por refuerzo, capacidades como el análisis de datos, la ejecución de código y la evaluación de calidad directamente al interior del modelo.

\begin{knowledgebox}{Por qué el análisis de datos empresarial es un buen banco de pruebas para el agentic training}
Este escenario cumple a la vez cuatro condiciones:
\begin{itemize}
  \item la cadena de la tarea es larga, no puede resolverse con una sola ronda de respuesta;
  \item la retroalimentación es relativamente verificable, lo que facilita construir el reward;
  \item las llamadas a herramientas son densas, lo que expone los problemas de action selection;
  \item el valor real es claro, y si la automatización tiene éxito es fácil de verificar desde el negocio.
\end{itemize}
\end{knowledgebox}

\subsection{El aprendizaje por refuerzo vuelve a ser clave en 2025}

En la discusión hay un consenso evidente: que el agentic training pudiera acelerarse en 2025 se debe en gran medida a que el RL por fin pasó de ser ``demasiado difícil de llevar a la práctica'' a ser un eslabón ``imposible de evitar''. Los invitados describieron la función del RL de manera muy pragmática: no es una fórmula mágica, sino la interfaz de entrenamiento que en verdad conecta environment, reward, tool use y la estrategia de largo plazo. Sobre todo en tareas de coding, GUI y tareas verificables, el reward sigue siendo costoso, pero ya puede formar un ciclo de optimización más cercano a la tarea real que el SFT puro.

\begin{warningbox}{El aprendizaje por refuerzo no es una llave mágica que resuelva al Agent de manera automática}
La postura del panel no fue ingenuamente optimista. Todos reconocieron que el reward, el environment, la eficiencia del entrenamiento y el espacio de exploración siguen siendo cuellos de botella. Es decir, el RL hace que el Agent se parezca más a un ``sistema entrenable'', pero no elimina la dificultad del diseño del sistema ni de la construcción de datos.
\end{warningbox}

\subsection{Resumen de la sección}

``Entrenar el modelo para que sea un Agent'' es el primer eje central de esta discusión. Su punto no está en agregarle al modelo un nuevo benchmark, sino en que el modelo adquiera de manera gradual la capacidad nativa de manejar entornos abiertos, tareas de cadena larga y retroalimentación de herramientas. El aprendizaje por refuerzo, las tareas verificables y los escenarios de trabajo reales son la razón por la que esta ruta en verdad se volvió viable en 2025.
